\section{Related Work}
\label{sec:related_work}

The literature relevant to this paper spans three partially connected threads: standards-oriented dynamic tests, estimator-family design papers, and low-inertia or IBR-focused disturbance studies. The gap is not a lack of candidate algorithms. The gap is a lack of comparative studies that preserve disturbance class and operating objective at the same time.

Standards and many benchmark papers still organize evidence around isolated compliance-style disturbances. IEEE C37.118.1 and IEC/IEEE 60255-118-1 emphasize ramps, magnitude steps, modulation tests, and limited phase events~\cite{IEEE_C37_118,IEC_IEEE_60255_118_1_2018}. That logic is indispensable because it defines minimum common tests, but it also privileges clean cases in which many estimators tie on protection-oriented risk. Wright \textit{et al.} show that even simple large phase steps can produce misleading RoCoF signatures~\cite{Wright2019}. Once such ambiguities appear, pass-fail compliance is no longer an adequate description of estimator behavior.

The estimator literature is richer on algorithms than on cross-family comparison. Window-based methods such as IpDFT, TFT, Prony, and ESPRIT trade spectral selectivity against observation delay and computational burden~\cite{Belega2014,PlatasGarza2011}. Loop-based methods remain attractive because they are lightweight and interpretable, but their transient behavior depends strongly on gain selection, bandwidth choice, and re-lock dynamics~\cite{Golestan2017,Wright2019}. Kalman-family estimators can outperform simpler methods when the state model matches the event, yet that advantage is conditional on model adequacy and noise handling~\cite{Ferrero2020EKF,Chamorro2021}. Reduced-order and data-driven methods introduce additional cost-accuracy trade-offs without, by themselves, resolving how the benchmark should be interpreted~\cite{Netto2018KoopmanKF,Ting2022}.

Recent low-inertia and IBR-focused studies widen that benchmark-design problem rather than close it. Converter-dominated systems change the time scale, disturbance composition, and recovery behavior of observable frequency dynamics~\cite{Milano2018LowInertia,Pinheiro2025,Ponce2025}. Hybrid, grid-forming, and interaction-rich studies emphasize precisely those non-ideal behaviors, yet they rarely compare a broad estimator registry over standard PMU tests, non-ideal waveform stressors, composite IBR transients, relay-oriented risk metrics, and computational cost inside one coherent reporting frame~\cite{Bernal2025,Zoghby2024}. Four needs follow directly from that literature: event-archetype scenarios rather than only isolated textbook tests, explicit non-ideal waveform stressors, joint latency-accuracy-risk-cost reporting, and comparative summaries that remain meaningful for deployment rather than only for algorithm taxonomies.

This manuscript addresses those needs by treating benchmark structure as a scientific contribution rather than as incidental experimental detail. The paper does not claim that one metric or one scenario family is sufficient. Instead, it argues that a journal-quality benchmark must preserve three objects simultaneously: disturbance family, metric definition, and deployment objective. A benchmark that compresses these into one scalar rank may be easy to read, but it removes the design decision the reader actually needs to make.
